\section{Formal verification}
\label{sec:formalization}

The Lean development accompanies the proofs above and records which parts
have been checked on their original mathematical objects. Its source and
verification evidence are available in the repository
\cite{repository2026}. The verified source snapshot is
\href{https://github.com/veridiscoverLab/fermionic-coherent-lean/tree/3164a3243fcc6cd675fa8a7493282a1f650a16d8}{\texttt{3164a3}};
the English manuscript was added subsequently without changing those proof
sources.

\subsection{End-to-end exterior-power theorems}

The formalization of Theorem~\ref{ext:convex-order} starts with
$\Lambda^p\C^n$, its actual unitary action, density operators, and normalized
Haar measure. It constructs creation and contraction, the conditional
density on the orthogonal complement, the double-Haar identity, and the
common occupation mixture. The induction and the strict equality argument
are proved from these constructions. In particular, the statement covers
every $0\le p\le n$, all mixed inputs, and every continuous convex test on
$[0,1]$; a strictly convex test has equality exactly at pure Slater
densities. No conditional-measure identity or full-density rigidity
statement is assumed as a separate hypothesis.

The formal Slater measure is the pushforward to the actual orbit of
rank-one Slater projectors. Its invariance, independence of the reference,
and agreement with group integration are proved. The entropy results then
use this same measurement. They include the sector dimension $\binom np$,
the normalized first moment, all positive real moments, the Wehrl minimum,
and the R\'enyi--Wehrl minimum for every real $q>0$, $q\ne1$, together with
the same equality characterization. These minimizer theorems are end to
end. Their coherent comparison values are defined by integrals; the
Beta-product evaluations, explicit Gamma products, harmonic-number
constants, and the limit as $q\to1$ have not been formalized.

The principal declarations are listed below. Each namespace has the prefix
\texttt{Fermionic.}; the repository prints the full types of all declarations,
including their quantifiers and domain assumptions.
\begin{center}
\small
\begin{tabular}{@{}p{0.28\linewidth}p{0.66\linewidth}@{}}
\toprule
Namespace & Main declarations \\
\midrule
\texttt{HusimiInduction} &
\texttt{exterior\_husimi\_convex\_order}\newline
\texttt{exterior\_husimi\_equality\_iff} \\[3pt]
\texttt{SlaterMeasure} &
\texttt{orbit\_husimi\_convex\_order}\newline
\texttt{orbit\_husimi\_equality\_iff} \\[3pt]
\texttt{ExteriorEntropy} &
\texttt{wehrl\_minimum}, \texttt{wehrl\_eq\_iff\_slater}\newline
\texttt{renyiWehrl\_minimum}\newline
\texttt{renyiWehrl\_eq\_iff\_slater} \\
\bottomrule
\end{tabular}
\end{center}

\subsection{The Gaussian-rank boundary}

For the actual even Spin-vacuum-orbit Gaussian dictionary on eight modes,
the formalized results give, for the original occupation vector $\Psi$
defined in Section~\ref{sec:gaussian-rank},
\[
 2\le\chi_G(\Psi)\le4.
\]
They include the actual Clifford representation, the four-term Gaussian
decomposition, purity of Gaussian orbit vectors, the vacuum-chart
annihilator calculation, and nonpurity of the target. Formal scale
invariance gives the same interval for $\Psi/2$. Its occupation-space
identification with $M^{\otimes2}$ is made in the paper; a separate
Hilbert tensor-product identification is not part of the verified roots.
Separately, the full
normal-family certificate in Section~\ref{sec:gaussian-rank} is checked for
arbitrary complex parameters, including zero parameters. The encoded
target has the nonzero certificate value $18063360$.

These results do not yet constitute an end-to-end formalization of
$\chi_G(M^{\otimes2})=4$. The remaining Lean work is the original
Fock-to-coordinate identification, covariance of the certificate under the
actual common transformations, and the reduction of every three-term
pure-spinor sum, including its degenerate cases. Section~\ref{sec:gaussian-rank}
gives the mathematical argument for these steps. The half-spin and odd-spin
convex-order theorems, their entropy consequences, and the applications in
Section~\ref{app:section} likewise have paper proofs here, rather than
completed end-to-end Lean proofs. The repository's
\href{https://github.com/veridiscoverLab/fermionic-coherent-lean/blob/main/TODO.md}{\texttt{TODO.md}}
records the outstanding formal statements and their paper-level derivations.

\subsection{Verification procedure and foundations}

The recorded runs use Lean~4.29.0 and mathlib at revision
\texttt{8a178386ffc0f5fef0b77738bb5449d50efeea95}.
All Lean and Lake execution took place on a CAB server through SSH.
Three suites passed clean project compilation with warnings treated as
errors, inspection of all selected declaration types and their transitive
axioms, and a fresh kernel replay of the imported proof environment.
The exterior suite contains 20 modules and 309 theorems; the rank
foundation and certificate suite contains 33 modules and 257 theorems;
the vacuum-chart suite contains 12 modules and 150 theorems. These suites
overlap and were verified in three separate runs. Their shared source
files are identical.

The proofs use no \texttt{sorry}, \texttt{admit}, added mathematical axiom,
or \texttt{native\_decide}. Finite certificate leaves use ordinary
kernel-checked decision procedures, and parameter identities use
kernel-checked algebraic proofs. The disclosed logical axioms are
\texttt{propext}, \texttt{Classical.choice}, and \texttt{Quot.sound}.
Fresh replay starts from an empty environment and rechecks the replayable
safe, non-partial constants with the same Lean kernel; axioms remain axioms.
It is not an independent implementation of the kernel. The dependency
versions are pinned, but their source files were not all recompiled from
scratch in each run. Source hashes, complete logs, and reproduction
instructions are included in the repository.
